\section{Overview and Objectives\label{sec:intro}}

The \textbf{objective} of this proposal is to develop a scalable Bayesian methodology to solve the inverse problem of reconstructing cardiovascular hemodynamic fields and cardiac structure from advanced medical imaging modalities.

Cardiovascular diseases are the leading cause of mortality in the United States.
Heart disease accounts for one-fifth of adult deaths \cite{heartstat}, with congenital heart conditions afflicting 1\% of newborns, contributing to 4.2\% of neonatal fatalities \cite{CHDstat1, petrini2010racial}.
Vital diagnostic information for heart disease lies in cardiovascular flow, which can be extracted from advanced medical imaging techniques like phase-contrast magnetic resonance imaging (PC-MRI) \cite{Nayak2015_PCMRI}, 4D Flow MRI \cite{maroun2023_4D}, and Color Doppler echocardiography (henceforth referred as Echo) \cite{brett2020_echo}.
These methods directly measure velocity, albeit with inherent noise and spatial-temporal averaging.
Computational fluid dynamics (CFD) \cite{morris2016_CFD} can reconstruct finer velocity details, pressure and wall shear stress.
Researchers have explored mapping medical images to flow fields through regression models trained on pairs of CFD and synthetic images~\cite{Ferdien2020_4DFlowNN,Rutkowski2021_4DFlowNN}.
However, CFD results are sensitive to variable vessel geometries, boundary/initial conditions, and differing fluid and material parameters across patients.

Quantifying uncertainty in hemodynamic field reconstructions is essential for clinical decision-making~\cite{Ninos2018_HemoUQ}. 
Understanding the uncertainty in reconstructed velocity fields can provide valuable information to clinicians to gauge the uncertainty in pressure drop across the stenosis, aiding treatment decisions.
Availability of uncertainty can also help optimize imaging protocols for maximum information gain.
Uncertainty quantification requires formulating data assimilation and calibration problems in a Bayesian way~\cite{Arzani2022_CardioML}.
While some attempts at Bayesian reconstruction of velocity fields exist \cite{Nguyen2022_HemoBayes, croci2021fast}, none enforce the underlying physics for joint field reconstruction.

Obstacles in the uncertainty quantification of hemodynamic fields stem from gaps in prior uncertainty quantification for vascular geometry, boundary conditions, and material parameters. 
Additionally, hemodynamic fields exhibit high dimensionality, rendering traditional Bayesian methods intractable \cite{gao2010composite}, especially when likelihood evaluations require calling a CFD solver~\cite{liu2021uncertainty}.
Approaches substituting CFD with regression models are impractical in high dimensions~\cite{Arzani2022_CardioML}.
Neural operators have made progress but require substantial training data or iterations when constrained by physics~\cite{li2020fourier, maul2023transient}.
Lastly, addressing model form errors remains underexplored.

The proposed methodology builds upon our preliminary work on information field theory (IFT)~\cite{alberts2023physics, hao2023information}, a theory rooted in Bayesian statistics over function spaces~\cite{ensslin2009information}, inspired by statistical and quantum field theories~\cite{parisi1988statistical, lancaster2014quantum}. 
It utilizes Feynman path integrals to define a prior probability measure over a function space, encoding beliefs about physical fields.
A physics-based measurement model connects the prior measure to data, conditioning it on observations.
Our contributions to IFT encompass encoding known physics (partial differential equations (PDEs)) in the prior measure without requiring a physical solver, addressing time evolution, and characterizing the joint posterior measure over fields and parameters via sampling and variational inference.

IFT is the most suitable framework for formulating the cardiovascular hemodynamics reconstruction problem, as it can fuse noisy data from multiple modalities, exploit known physics without relying on a PDE solver, handle ill-posed problems, and identify model form errors.
Consequently, IFT can characterize the joint posterior measure over fields and parameters, quantify uncertainty in reconstructed fields, improve medical image analysis, and enhance clinical decision-making.

We will accomplish our objective via the following specific aims (as shown in Fig. \ref{fig:overview_of_IFT}):
\begin{itemize}[leftmargin=*]
\item \textbf{Specific Aim 1: Theoretical formulation of the reconstruction problem.} The objective is to formulate the inverse problem of identifying the hemodynamic fields of interest within the IFT framework. 
This involves developing physics-informed prior probability measures for hemodynamic fields (\emph{Task 1.1}), physics-based PC-MRI, 4D Flow MRI, and Echo measurement operators (\emph{Task 1.2}), and a unifying approach to solve the co-registration problem (\emph{Task 1.3}).

\item \textbf{Specific Aim 2: Parameterization of time-evolving fields.}  The objective is to study the properties of different parameterizations of spatiotemporal hemodynamic fields.
This involves parameterizing spatial fields \emph{(Task 2.1)} and their time evolution \emph{(Task 2.2)}.

\item \textbf{Specific Aim 3: Numerical algorithms for characterizing joint posterior.} The objective is to develop scalable numerical schemes to characterize the joint posterior over hemodynamic fields and physical parameters.
This involves developing advanced stochastic gradient Langevin dynamics (SGLD) algorithms to sample from the joint posterior (\emph{Task 3.1}) and an amortized variational inference approach that learns the map from data to the joint posterior (\emph{Task 3.2}).

\item \textbf{Specific Aim 4: Verification, validation, and demonstration.} The objective is to verify, validate, and demonstrate our methodology on a series of benchmarks.
This involves verification using synthetic experimental data with known ground truth (\emph{Task 4.1}), validation using in vitro experimental data (\emph{Task 4.2}), and demonstration of the methodology on clinical data (\emph{Task 4.3}).

\end{itemize}

\begin{figure}%[htbp]
    \centering
    \includegraphics[width=\textwidth]{figs/overview-with-specific-aims.pdf}
    \caption{Schematic overview of the research plan.}
    \label{fig:overview_of_IFT}
\end{figure}

\section{Intellectual Merit\label{sec:im}}

The primary project outcome is a next-generation statistical methodology for quantifying uncertainty in hemodynamic field reconstructions from advanced imaging. Our team will create advanced mathematical models that use physical laws to inform prior probability measures over hemodynamic fields and the vascular structure. We will also develop probabilistic models for medical imaging techniques. To integrate data from these different sources, we will create a method for solving the problem of data co-registration. Our approach will include causality-respecting parameterizations of spatiotemporal fields, and we will assess their convergence rates. We plan to create scalable computational schemes for sampling and variational inference to characterize joint posteriors effectively. Additionally, we will implement these methodologies in open-source software optimized for GPU clusters. To verify and validate our approach, we will use synthetic examples and in vitro phantoms, and we will compare our methods against existing ones using clinical data.

\section{Review of Relevant Literature\label{sec:review}}

\subsection{Review of Cardiovascular Hemodynamics Reconstruction Methods}


Cardiovascular conditions, including coronary artery disease, heart failure, and stroke, are the predominant causes of death in the United States. Specifically, heart disease accounts for approximately one out of every five adult deaths \cite{heartstat}. Congenital heart defects, which occurs in 1\% of newborns, can necessitate life long care depending on its severity \cite{CHDstat1}. Early identification and diagnosis of cardiovascular diseases is critical in saving patient lives. Over the years, advancements in cardiac MRI and echocardiography imaging technologies continue to better sever this purpose.

PC-MRI~\cite{pelc1991_PCMRI} and 4D Flow MRI~\cite{markl2014_4d} are powerful clinical modalities for evaluating cardiovascular diseases owing to their ability to visualize and quantify complex flow structures within the heart chambers\cite{Nayak2015_PCMRI,maroun2023_4D}. Limitations related to long scan times and imaging artifacts have been addressed through recent advances in hardware~\cite{Stankovic2014_4Dreview}, accelerated imaging approaches~\cite{soulat2020_4Dreview} and pre-processing methods to correct artifacts (muti-venc aquisitions~\cite{Schnell2017_4Ddualvenc,ha2016_4Dmultivenc}, enforcing flow physics~\cite{zhang2021_4Dphaseunwrap}). However, the most pressing issue with MRI is its low spatial resolution ($\sim$ 1mm voxel size). This spatial averaging introduces uncertainty in hemodynamic evaluations, especially in regions with rapid changes in the velocity field. Multi-modality data assimilation methods employ CFD, particle image and tracking velocimetry (PIV/PTV) to address this issue \cite{bakhshinejad2017_4DmriCFD, zhang2022_4Dsparserep}. Recently, numerous machine learning-based techniques have emerged for super-resolution flow reconstructions \cite{Ferdien2020_4DFlowNN,Rutkowski2021_4DFlowNN,Arzani2022_CardioML}. The other major challenge with flow MRI is the long post-processing time, primarily due to manual interventions in image segmentation and co-registration \cite{zhuang2021_4Dreview}. Existing techniques for segmentation are often semi-automatic \cite{yushkevich2006_ITKSNAP}, requiring substantial user input, which introduces another layer of uncertainty and potential for error. Additionally, automatic techniques proposed suffer from being static segmentation techniques, making them inapt for cardiovascular applications~\cite{rothenberger2023_4DflowSDM}, or are based on neural networks, thus requiring huge training datasets for accurate patient-specific segmentation \cite{berhane2020_4Dseg,ziegler2021_MRIseg}.

Color Doppler Echocardiography is a widely used, quick and inexpensive, non-invasive diagnostic tool for quantitative blood flow assessment \cite{cohen1996_echo}.  The imaging is based on Doppler effect by which flow measurements recorded (as a colour map) only indicate the speed of flow towards or away from the transducer probe. This necessities accurate recovery of 2D/3D blood velocity fields from a set of 1D velocity components measured along the probe scan lines. Recent methods apply flow physics constraints such as mass conservation \cite{garcia2010_EchoIVFM,assi2017_EchoIVFM} and streamfunction-vorticity \cite{ brett2020_echo} for improved reconstructions.
%There is hardly any literature on CD-based automatic segmentation from echos \cite{jiang2021_echoseg,saad2006_echoseg}. Almost all the research is only focused on segmentation from B-mode images.

Despite all these advances, the need to account for imperfections in the medical imaging and reporting the effect of model uncertainties in the study is critical for clinical acceptability of any method. Added on, the ill-posed nature of the reconstruction problem, resulting from unknown boundary conditions and patient-specific parameters, adds more complications. An exhaustive summary of all efforts towards uncertainty quantification in cardiovascular flows is well presented in \cite{Arzani2022_CardioML} and \cite{Ninos2018_HemoUQ}. Amongst all the methods, data assimilation Bayesian inference methods are considered the most consistent approach to tackle the problems at hand.

\subsection{Review of Bayesian Methods for Data Assimilation}

Data assimilation blends theoretical dynamical models with experimental observations to identify the underlying dynamical state. 
The models often take the form of time-dependent PDEs.
Observations usually consist of noisy measurements, such as pressure or velocity at a specific location.
Multi-modal data assimilation finds utility in fields as diverse as numerical weather prediction~\cite{xue_advanced_2003}, geophysics~\cite{petra_computational_2014}, structural health monitoring~\cite{Chatzi2009, Maes2018, Yang2007, Chatzi2010}, and bioengineering~\cite{raissi_hidden_2020,rajendran_uncertainty_2020,zhang2019_4Dpressure,meyers_development_2018}.

Data assimilation uses Bayesian methods to determine the posterior distribution of the system states conditional on all observations.
For linear dynamics with Gaussian noise, the Kalman filter offers an efficient solution~\cite{Kalman1960}.
Its extended~\cite{Sorenson1966, Schmidt1981} and unscented~\cite{Julier1995, LUND2020106512} variants can handle non-linear cases but struggle with scalability with state dimensions.

Particle filters~\cite{Gordon1993, Doucet2001, Cappe2007} provide sampling-based approximations of the posterior, accommodating non-linear dynamics and non-Gaussian noise.
These filters scale more effectively than Kalman filter variants as state dimensionality increases. 
However, their applicability hinges on the computational efficiency of simulating the dynamical system. 
In cases where this is not feasible, an inexpensive surrogate model must suffice, which presents challenges in high-dimensional state spaces.

% Particle filters~\cite{Gordon1993, Doucet2001, Cappe2007} construct sampling approximations of the posterior over states and work with non-linear dynamics and non-Gaussian noise. 
% They scale better than variations of Kalman filters with increasing state dimensionality.
% If propagating the underlying dynamical system is computationally efficient, particle filters can work in real-time.
% Otherwise, one has to replace the physical simulator with an inexpensive surrogate.
% The latter is ubiquitously tricky if the state space is high-dimensional, as for time-dependent systems of non-linear PDEs.

Variational filters (VF)~\cite{Jordan1999, Blei2017} represent another approach that seeks an analytical approximation of the posterior within a predefined set of distributions, known as guides. 
This approximation minimizes the Kullback-Leibler divergence between the guide set and the true posterior. 
While effective for high-dimensional problems, variational filters can introduce bias if the guide set does not include the true posterior. 
Even though it has never been attempted before in a filtering context, it is also possible to use invertible DNNs and normalizing flows to construct sets of guides that come arbitrarily close to containing the true posterior~\cite{ardizzone2018analyzing}.
Another way to avoid systematic bias is to recast VF using Stein variational gradient descent~\cite{Liu2016, Liu2017}. 
This technique bypasses the need for a guide by directly constructing a particle approximation to the posterior.
Researchers have proposed several versions of VFs~\cite{Sun2012, Salimans2014, Kim2019, Broderick2013, Lund2020VariationalIF}. 
However, there is still a lot of research needed to assess which one, if any, performs consistently well and in which types of applications.
In any case, the real-time implementation of VFs requires a computationally inexpensive dynamical model.

% Variational filters (VF)~\cite{Jordan1999, Blei2017} is a new class of methods seeking an analytical approximation to the posterior over the states within a set of candidate distributions known as guides.
% One constructs this analytical approximation by minimizing the Kullback-Leibler divergence between the set of guides and the true posterior.
% With variational inference, one can deal with high dimensionality by suitably parameterizing the set of guides.
% Typically, variational filters have a systematic bias because the set of guides may not contain the true posterior.
% Even though it has never been attempted before in a filtering context, it is also possible to use invertible DNNs and normalizing flows to construct sets of guides that come arbitrarily close to containing the true posterior~\cite{ardizzone2018analyzing}.
% Another way to avoid systematic bias is to recast VF using Stein variational gradient descent~\cite{Liu2016, Liu2017}. 
% This technique bypasses the need for a guide by directly constructing a particle approximation to the posterior.
% %Stein variational gradient descent has also not been explored in a filtering context.
% Researchers have proposed several versions of VFs~\cite{Sun2012, Salimans2014, Kim2019, Broderick2013, Lund2020VariationalIF}. 
% However, there is still a lot of research needed to assess which one, if any, performs consistently well and in which types of applications.
% In any case, the real-time implementation of VFs requires a computationally inexpensive dynamical model.

Physics-informed neural networks (PINNs) \cite{raissi2019physics} is another prevalent approach in scientific machine learning.
The core of PINNs is defining state variables in the system using deep neural networks (DNNs) and ensuring these networks comply with relevant physical laws through additional PDE residual regularizers.
This structure promotes efficient learning with limited data.
However, standard PINNs do not account for uncertainty. 
%leading to critiques~\cite{chuang2022experience}.
Bayesian PINNs (B-PINNs)~\cite{yang2021b} emerged to incorporate uncertainty into PINNs.
A B-PINN consists of a Bayesian DNN, which puts a prior over the network parameters~\cite{mackay1992practical, neal2012bayesian}.
These networks then use methods like Hamiltonian Monte Carlo~\cite{betancourt2017conceptual} or variational inference ~\cite{blei2017variational} to estimate the posterior distribution of these parameters. 
The treatment of physics in B-PINNs can lead to training issues.
The physics is enforced via a fictitious Gaussian likelihood of zero PDE residuals on a predetermined set of collocation points.
It is unclear how to select or interpret the variance of this likelihood.
Further, this treatment prevents the physics in the entire domain from being enforced.
As the number of collocation points increases, the posterior will collapse to a single field, undesirable in model-form uncertainty or ill-posed problems.

% One of the recent most ubiquitous approaches to scientific machine learning is physics-informed neural networks (PINNs) \cite{raissi2019physics}.
% The idea behind PINNs is to parameterize state variables in the system as deep neural networks (DNNs) and force the network to respect any applicable physical laws using additional regularizers in the form of PDE residuals. 
% These regularizers lead to efficient learning under limited data.
% However, plain PINNs cannot capture uncertainty, and some groups have become critical in their use over other approaches~\cite{chuang2022experience}.
% Bayesian PINNs~\cite{yang2021b} were introduced as a way to quantify uncertainty using PINNs.
% A Bayesian PINN consists of a Bayesian neural network, which puts a prior over the network parameters~\cite{mackay1992practical, neal2012bayesian}.
% Then, Hamiltonian Monte Carlo~\cite{betancourt2017conceptual} or variational inference~\cite{blei2017variational} are used to estimate the posterior distribution over the network parameters.
% The treatment of physics in Bayesian PINNs can lead to training issues.
% The physics is enforced via a fictitious Gaussian likelihood of zero PDE residuals on a predetermined set of collocation points.
% It is unclear how to select or interpret the variance of this likelihood.
% Further, this treatment prevents the physics in the entire domain from being enforced.
% As the number of collocation points increases, the posterior will collapse to a single field, undesirable in model-form uncertainty or ill-posed problems.

%In summary, Bayesian methods for data assimilation have evolved to address various challenges but also present opportunities for further research, particularly in scalability, bias correction, and the integration of physical laws.


\subsection{Remaining Knowledge Gaps}
The proposed work aims to address and overcome the following major drawbacks:
\begin{itemize}[nosep,leftmargin=*]
\item Uncertainty quantification in hemodynamic fields is not well understood (\emph{Specific Aim 1});
\item Curse of dimensionality regarding the representation of hemodynamic fields (\emph{Specific Aim 2});
\item Reliance on computationally expensive CFD solvers (\emph{Specific Aims 1\&3});
\item Lack of scalable algorithms for batch Bayesian data assimilation (\emph{Specific Aim 3}).
\end{itemize}

\section{Background on Information Field Theory and Preliminary Work\label{sec:prelim}}

\subsection{Brief introduction to information field theory}
\label{sec:IFTreview}
IFT is a framework for Bayesian inference over fields.
A field is a physical quantity $\phi$ with a value for every point in space $x$ and time $t$.
One can have scalar, vector, or tensor fields.

\textbf{Information Hamiltonian: Priors over fields encode our beliefs about them.} IFT begins by defining a prior probability measure over fields, which we formally denote by $p(\phi)$.
% IB says: Commenting it out because I want to see ho much space it takes.
%\begin{figure}[htbp]
%    \centering
%    \includegraphics[width=1\textwidth]{figs/randomfields.pdf}
%    \caption{probability measures over reals vs fields}
%    \label{fig:randomfields}
%\end{figure}
This measure is defined through the \textit{information Hamiltonian}, $H(\phi)$, which is the negative logarithm of $p(\phi)$ up to an additive constant, $H(\phi) = -\log p(\phi) + \text{const}.$
The information Hamiltonian is a functional of the field that typically includes a spatiotemporal integral of a combination of field expressions. 
For example, the quadratic form $H(\phi) = \int dx\: dx'\: \phi(x)S(x,x')\phi(x')$ defines a zero-mean Gaussian process (GP) with covariance $S^{-1}(x,x')$, the Green's function of the local operator $S(x,x')$ \cite{ensslin2019information}.
We can use the information Hamiltonian to encode our prior beliefs about the fields, e.g., regularity, scale, periodicity, interactions between fields, and more.

\textbf{Data Hamiltonian: Measurement model connects fields to data.} IFT connects the fields to the observed data $d$ using a probabilistic model of the measurement process, $p(d|\phi)$ (likelihood).
The likelihood depends on a \emph{measurement operator} $R$, e.g., $p(d|\phi)=p(d|R\phi)$.
The expression $R\phi$ is what we expect to observe if the underlying field is $\phi$.
For example, for a point measurement at location $x'$, $R\phi = \int dx \: \delta(x-x')\phi(x)$.
We work with the \emph{data Hamiltonian}, $H(d|R\phi) = -\log p(d|\phi) + \text{const}$.
For example, the data Hamiltonian $H(d|R\phi) = \frac{1}{2}(d-R\phi)^TN^{-1}(d-R\phi)$ corresponds to a Gaussian measurement model with mean $R\phi$ and covariance matrix $N$.

\textbf{Posterior information Hamiltonian: What we know after seeing the data.}
One can find the posterior over fields using Bayes's rule. The \emph{posterior Hamiltonian} is:
\begin{equation}
\small
\label{posterior-hamiltonian}
\underbrace{H(\phi|d)}_{\text{posterior information Hamiltonian}} 
= \underbrace{H(d|R\phi)}_{\text{data Hamiltonian (measurement operator}\;R)} 
+ 
\underbrace{H(\phi).}_{\text{prior information Hamiltonian}}
\end{equation}
Under special cases, Eq.~(\ref{posterior-hamiltonian}) yields the posterior of a GP.
In general, the result is a non-Gaussian random field.
The posterior encodes all uncertainty about the fields after seeing the data.

\textbf{The partition function complicates things.}
Suppose the information Hamiltonian depends on some parameters $\lambda$, say $H(\phi|d) = H_\lambda(\phi|d)$.
Formally, the posterior measure is:
\[
\small
\underbrace{p(\phi|d,\lambda) = \frac{\exp\left\{-H_\lambda(\phi|d)\right\}}{Z(\lambda)},}_{\text{posterior probability measure over fields}}
\;\text{where}\;
\underbrace{Z(\lambda) = \int \mathcal{D}\phi\: \exp\left\{-H_\lambda(\phi|d)\right\}.}_{\text{the normalization-constant function needs a functional integral}}
%p(\phi|d) = \frac{\exp\left\{-H(\phi|d)\right\}}{Z},\;\text{where}\;Z = \int \mathcal{D}\phi \exp\left\{-H(\phi|d)\right\}.
\]
The normalization constant, $Z(\lambda)$, is called the \emph{partition function}, a function of all the Hamiltonian parameters.
The corresponding integral is a \emph{functional integral}.
It is not a mathematically well-defined quantity, e.g., it is infinite even for GPs.
In practice, physicists tame these infinities using renormalization groups \cite{lancaster2014quantum}, effectively restricting the smallest possible scale a field can represent.
%Employing a parameterization of the fields is equivalent.
When characterizing the fields conditional on the Hamiltonian parameters, one can ignore the partition function.
However, the partition function cannot be dropped if one wants to infer the parameters $\lambda$.
For this reason, the partition function is the biggest roadblock to understanding and implementing IFT.
Fortunately, in~\cite{alberts2023physics} we showed that one only needs the derivatives of the partition function with respect to $\lambda$, which can be approximated by Monte Carlo.


\subsection{Our contributions to physics-informed, information field theory}

We started constructing probability measures that encode known physics in 2020, inspired by our previous work on variational principles for solving stochastic PDEs~\cite{karumuri_simulator-free_2020}, reaching back to 2016~\cite{bilionis2016probabilistic}.
We presented preliminary results at the SIAM CSE conference of 2021~\cite{AlbertsBilionis2021}.
However, we could continue our program only after studying the mathematics of statistical and quantum field theories and discovering the literature on IFT.
We presented updated results at the SIAM UQ conference of 2022~\cite{AlbertsBilionis2022}~and managed to publish our first paper~\cite{alberts2023physics} in 2023, which was presented at UNCECOMP 2023~\cite{AlbertsBilionis2023}.
In the under review paper~\cite{hao2023information}, we extended our methodology in various ways.
In what follows, we summarize our contributions while providing evidence that IFT is the most appropriate theory for formulating the cardiovascular
hemodynamics reconstruction problem.

\textbf{IFT can fuse noisy data from multiple modalities.}
We can construct a likelihood from multiple measurement modalities using the data Hamiltonian, exploiting the fact that the data Hamiltonian is additive if the experiments are independent conditional on the fields.
To illustrate the idea, separate the data vector according to the $n$ different measurement modalities: $d = (d_1,\dots, d_n)$.
Then, one needs a different measurement operator $R_i$ and a corresponding data Hamiltonian $H_i$ for each modality.
The data Hamiltonian for the combined dataset is:
\begin{equation}
\small
    \underbrace{H(d|R\phi)}_{\text{data Hamiltonian}} = \underbrace{H_1(d_1|R_1\phi)+\dots + H_n(d_n|R_n\phi).}_{\text{additive decomposition over (independent) experimental modalities}}
    \label{eqn:data_hamiltonian}
\end{equation}

\textbf{IFT can exploit known physics without needing a PDE solver.}
Typically, the physics is given by a PDE, and we derive an information Hamiltonian, $H(\phi)$, the minimization of which yields a PDE solution.
In this way, the corresponding prior probability measure puts more weight on fields with a smaller $H(\phi)$, and it is peaked at PDE solutions functions.
$H(\phi)$ is proportional to field energy if a variational principle exists; otherwise, one can use the integrated square residual of the PDE.
To illustrate, let $\phi$ be a scalar spatial field that satisfies the Poisson equation:
\begin{equation}
\small
\label{eqn:prior_hamiltonian}
\underbrace{\nabla\cdot \left(a\nabla \phi\right) + f = 0}_{\text{field PDE}}\Rightarrow
\underbrace{H(\phi) = \beta \int dx 
\left(\frac{a}{2}|\nabla\phi|^2 - \phi f\right)}_{\text{energy-based Hamiltonian}}\text{or}
\underbrace{H'(\phi) = \beta'\int dx \left[\nabla\cdot \left(a\nabla \phi\right) + f\right]^2.}_{\text{integrated square residual Hamiltonian}}
\end{equation}
The $\beta$ parameter is positive and ensures that the information Hamiltonian is unitless.
In our publications, we demonstrated analytically and numerically that $\beta$ is inversely proportional to the field variance and that the measure collapses to PDE solutions as $\beta$ tends to infinity.
We have also developed approaches that can characterize the field posterior (either through sampling or variational inference) using only Monte Carlo estimates of the information Hamiltonian integrals.
%integrals that appear in the information Hamiltonian.
%Therefore, we can interpret $\beta$ as the confidence in the correctness of the physics, a point we analyze more in a subsequent paragraph.
%The prior is defined in such a way that fields which are closer to minimizing the energy are a priori more likely.
%Dealing with the physics-informed prior requires evaluating the integral contained in $U$, which can be evaluated efficiently by casting the integral as an expectation and sampling space.
%We have developed an algorithm based on stochastic gradient Langevin dynamics \cite{welling2011bayesian}, where one spatial point is sufficient to sample the prior.

\textbf{IFT can deal with ill-posed problems.}
Inverse problems may have multiple possible solutions, commonly due to insufficient data, e.g., missing boundary conditions and unobserved fields.
IFT deals with such ill-posed problems in three ways.
First, as a Bayesian method, it yields a probability measure over the space of solutions and not a point solution, capturing multiple possible solutions to the problem.
This probability measure encodes the epistemic uncertainty induced by limited data.
Second, it can enforce prior beliefs about the field's regularity and smoothness.
Third, and this constitutes our unique contribution to the core theory, it rules out fields unlikely to satisfy physics.
The second two points, unique to IFT, narrow the possibilities and make exploring the posterior numerically easier.
In Fig.~\ref{fig:3bcompiled}, we present an example from our paper~\cite{alberts2023physics}.
The task is to infer all parameters and the source term of the non-linear PDE $D\frac{d^2\phi}{dx^2}-\kappa\phi^3=f$, using only field measurements.
\begin{figure}[htbp]
    \centering
    \begin{subfigure}[b]{0.32\textwidth}
        \centering
        \includegraphics[width=\textwidth]{figs/example03b_theta_post.pdf}
        \caption{Posterior over parameters.}
        \label{fig:ex3b_post_params}
    \end{subfigure}
    \begin{subfigure}[b]{0.32\textwidth}
        \centering
        \includegraphics[width=\textwidth]{figs/example03b_source.pdf}
        \caption{Posterior of unobserved source.}
        \label{fig:ex3b_source}
    \end{subfigure}
    %\begin{subfigure}[b]{0.475\textwidth}
    %    \centering
    %    \includegraphics[width=\textwidth]{figs/example03b_fitted_prior_predictive.pdf}
    %    \caption{Fitted prior predictive.}
    %    \label{fig:ex3b_prior}
    %\end{subfigure}
    \begin{subfigure}[b]{0.33\textwidth}
        \centering
        \includegraphics[width=\textwidth]{figs/example03b_fitted_post_predictive.pdf}
        \caption{Posterior of field.}
        \label{fig:ex3b_field_post}
    \end{subfigure}
    \caption{Ill-posed problem example solved using IFT. Reproduced from~\cite{alberts2023physics}.}
    \label{fig:3bcompiled}
\end{figure}

\begin{wrapfigure}{r}{0.5\textwidth}
  \begin{center}
    \includegraphics[trim={0 0.3cm 0 0.1cm},clip,width=0.48\textwidth]{figs/example02a.pdf}
  \end{center}
  \caption{Detecting model-form error using IFT. Reproduced from~\cite{alberts2023physics}.}
  \label{fig:ex2compiled}
\end{wrapfigure}
\textbf{IFT can detect model-form errors.}
When implementing models that employ physical knowledge, it becomes important to quantify the uncertainty that arises from the selection of the physics.
Model-form uncertainty is the epistemic uncertainty from an imperfect physical model~\cite{draper1995assessment, kennedy2001bayesian}.
The recent surge in popularity of physics-informed methods has led to a new generation of methods for quantifying model-form uncertainty~\cite{sargsyan2019embedded, park2014bayesian, riley2011quantification, geneva2019quantifying, brynjarsdottir2014learning}.
The available methods typically couple a solver for the physics with a technique from uncertainty quantification.
This involves resolving the problem multiple times, creating a bottleneck.
The IFT approach we propose avoids this issue by employing the $\beta$ parameter of Eq.~(\ref{eqn:prior_hamiltonian}).
As mentioned earlier, we have shown that the inverse of $\beta$ controls the prior field variance.
If $\beta$ tends to infinity, the prior field variance collapses to the PDE solution, and we have complete trust in the physics.
If $\beta$ is zero, the prior measure becomes flat, and we do not believe the physics to be correct.
We automatically quantify model-form error by treating $\beta$ as a hyperparameter and inferring it along with the field.
In Fig.~\ref{fig:ex2compiled}, we report the results of synthetic numerical experiments where the posterior of $\beta$ is tracked against controllable model correctness.

\textbf{IFT can represent dynamics.}
The typical approach in IFT for dynamical systems is to employ information field dynamics (IFD)~\cite{ensslin2013information}.
The general idea of IFD resembles KF's for fields.
The two issues with this approach are that it treats time as discrete and it requires that the priors and posteriors remain sufficiently Gaussian.
In our applications, the use of physics-informed priors practically guarantees that the priors will be non-Gaussian, especially in the case of ill-posed inverse problems.
Thus, new techniques to handle time evolution are required.

In~\cite{hao2023information} we developed a new method for representing dynamics through IFT.
We did it for finite dimensional dynamical systems.
Suppose that the dynamical system's state is $\mathbf{y}$ and that it satisfies an ODE.
We derived an information Hamiltonian based on the integrated square residual and we proved that it corresponds to the probability law of a stochastic ODE in the It\^o sense:
$$
\small
\underbrace{\dot{\mathbf{y}} = \mathbf{f}(\mathbf{y},t)}_{\text{ODE}}\Rightarrow
\underbrace{H(\mathbf{y}) = \beta\int_0^Tdt|\dot{\mathbf{y}} - \mathbf{f}(\mathbf{y},t)|^2}_{\text{information Hamiltonian}}\Leftrightarrow
\underbrace{d\mathbf{Y}_t = \mathbf{f}(\mathbf{Y}_t,t)dt + (2\beta)^{-\frac{1}{2}}d\mathbf{B}_t.}_{\text{stochastic differential equation}}
$$
Here $\mathbf{B}_t$ is a multi-dimensional Wiener process.
We see here that $\beta$ controls the process noise.
However, there is difficulty in finding an equivalent treatment for time-dependent stochastic PDEs.
We make a specific proposal in \emph{Task 2.2}.

\begin{wrapfigure}{r}{0.5\textwidth}
  \begin{center}
    \includegraphics[trim={0 0.3cm 0 0.3cm},clip,width=0.48\textwidth]{figs/posterior-inference-distributed.pdf}
  \end{center}
  \caption{IFT naturally lends itself to parallelization schemes that allow inference to scale up to large computational domains and datasets.}
  \label{fig:posterior-inference-distributed}
\end{wrapfigure}
\textbf{IFT posterior inference can be scaled.}
We have shown that inference, either via sampling~\cite{alberts2023physics} or variational inference~\cite{hao2023information}, can proceed without evaluating the Hamiltonians exactly.
We developed algorithms that converged to the posterior using only noisy sampling estimates of the gradients of the Hamiltonians.
For the spatiotemporal integrals, one needs the ability to sample uniformly in the computational domain.
For the data Hamiltonian, one needs to process batches of the observed data.
The significance of this observation is twofold.
First, we get iterative algorithms with a fixed computational and memory cost per iteration, irrespective of the computational domain size or the amount of data we need to process.
%This is a hallmark of IFT, as the continuous framework removes dependencies on any grid-specific properties~\cite{ensslin2013information}.
Second, the algorithms can exploit distributed computing resources using a data parallelism framework (Fig.~\ref{fig:posterior-inference-distributed}).
The idea is to copy the physical fields in all available computational tasks but split the data among the tasks.
In our case, one would distribute amongst processes the computational domain points (for evaluating the information Hamiltonian) and the observations (for evaluating the data Hamiltonian).
This allows scaling to data sets and computational domains that do not fit into the memory of any single task, and it enables the parallel evaluation of the necessary gradients that drive the inference algorithms.
Finally, one needs to do a reduction to aggregate the sampling averages to move to the next step.
%These parallelization schemes, and more, will be further studied in \emph{Specific Aim 3}.


\section{Research Plan\label{sec:research-plan}}

\subsection*{Specific Aim 1: Theoretical formulation of the reconstruction problem.}
\paragraph{Introduction.}
The objective is to formulate the inverse problem of identifying the hemodynamic fields of interest within the IFT framework. We will merge physical knowledge with data coming from advanced imaging modalities. This requires specifying a physics-informed prior probability measure over fields and the likelihood which describes the measurement process. The prior will contain hyperparameters or hyperfields, including hemodynamic fields that are not directly observed. The physics is given by PDEs (Eikonal, Continuity, Navier-Stokes), and we will investigate the definition of prior probabilities that encode them. The PDEs may be coupled with boundary/initial conditions, which can also be encoded in the prior. For the likelihood, we must derive the measurement operators, i.e., construct physical models that connect the hemodynamic fields to the medical modality images.
Finally, we need to address the co-registration problem when we fuse data from multiple modalities.

\textbf{Task 1.1: Develop physics-informed prior probability measures for hemodynamics.}
The fields of interest are the geometry of the vessel, the velocity and pressure of blood inside the vessel, and the strain and stress tensor over the vessel.
We are not going to attempt an explicit reconstruction of vessel strain and stress, as they require modeling expertise beyond our teams's.

We will assume that the vessel is a time-evolving, 2D smooth surface.
We will use a signed distance function (SDF), $g_0$, to represent this surface at the initial timestep.
The SDF $g_0(x)$ gives the minimum distance of $x$ from the vessel, and it is negative if $x$ is inside the vessel and positive otherwise.
The vessel corresponds to the zero contour of $g_0$.
The SDF satisfies the Eikonal equation, $|\nabla g_0| = 1$.
The SDF that represents the geometry at any time $t$, satisfies the level set equation, $\frac{\partial g}{\partial t} + \mathbf{v}_g^T\nabla g = 0$, with initial condition $g(x,0) = g_0(x)$.
Here, $\mathbf{v}_g$ is an unknown spatiotemporal vector field determining how the surface moves.
Using this information, we will write down a joint probability measure that factorizes as $p(g, g_0|\mathbf{v}_g) = p(g|g_0, \mathbf{v}_g) p(g_0)$ and the objective becomes to find information Hamiltonians for each term.
There is always a default starting point to construct such information Hamiltonians since we can use the PDE's integrated squared residual (ISR) that the fields satisfy.
For example, the information Hamiltonian for the initial SDF could be:
$$
\small
H(g_0) = \beta_{g_0}\int dx \left(|\nabla g_0| - 1\right)^2,
$$
where $\beta_{g_0}$ is a positive scaling coefficient to be determined.
We will investigate whether variational principles may be better choices than the ISR.
To complete the measure, one must specify a prior over the unknown field $\mathbf{v}_g$.
The default choice would be to pick a GP prior.

The velocity vector field, $\mathbf{v}$, and the pressure scalar field, $\tau$, are only defined inside the vessel, i.e., wherever $g < 0$, and they satisfy the Navier-Stokes (NS) equations.
We need to enforce no-slip boundary conditions, which means that $\mathbf{v} = 0$ wherever $g=0$, which we plan to do through explicit parameterization, see \emph{Task 2.1}.
Suppose we denote by $\mathbf{v}_0$ and $\tau_0$ the initial conditions of velocity and pressure, respectively. In that case, we will encode the information above in a prior probability measure of the form $p(\mathbf{v}, \tau|g, \mathbf{v}_0, \tau_0)$.
Again, the starting point will be to construct an information Hamiltonian using the ISR of the NS, but we will explore variational formulations of the NS~\cite{finlayson1972existence, vecile1984extended, bramble1994variational, manouzi2005mixed, taha2023variational, bistafa2023lagrangians}.
The no-slip boundary conditions can be enforced parametrically, see \emph{Task 2.1}.

\begin{figure}%[htbp]
    \centering
    \includegraphics[width=\textwidth]{figs/MRI_measurement_process.pdf}
    \caption{Schematic of 4D Flow MRI measurement process. Images are adopted from \cite{van1999_MRIbasics,markl2014_4d,mriquestions_kspace}.}
    \label{fig:MRI_measurement_process}
\end{figure}

\subsubsection*{Task 1.2: Develop physics-based measurement operators.}
We will construct mapping functions that model the measurement process of PC-MRI, 4D Flow MRI, and Echo. In general terms, we formulate the \emph{reference} measurement operator for each imaging modality.
We use the term ``reference'' because we operate at a single time $t$ and under the assumption that the geometry is perfectly registered.
Putting all timesteps and modalities together requires addressing the co-registration problem, which we do in \emph{Task 1.3}.

%Use $\boldsymbol{\phi}$ to denote the tuple of all relevant fields to be reconstructed, i.e., $\boldsymbol{\phi}=(g,\mathbf{v},\tau,\mathbf{v}_g)$.
%The reference operator for modality $i$ at time $t$ is denoted by $R_{i,t}^\text{ref}$.
%In the case of 4D Flow MRI, for example, this operator takes the tuple of field $\boldsymbol{\phi}$ and it produces a 3D array (potentially vector valued).
%If we let $x_{jkr}$ be the central coordinate of a voxel in this 3D image.
%Then, we can represent the $jkr$ component of the operator as:
%$$
%\small
%R_{i,t,jkr}^{\text{ref}}\boldsymbol{\phi} = \int dx' dt' r_i^{\text{ref}}(x_{jkr}-x',t-t')\mathbf{v}(x',t'),
%$$
%where $r_i^{\text{ref}}$ is a reference measurement kernel.
%Notice the assumption that the kernel is stationary.

PC-MRI and 4D Flow are based on phase-contrast imaging, where the blood velocity information is encoded in the phase of the complex-valued MRI signal. In principle, the measurement operator would entail (Fig. \ref{fig:MRI_measurement_process}) solving the Bloch equations for the total magnetic moment from all spins in the sample (step 1-2), modeling the k-space acquisition (step 2-3), reconstructing the MRI (magnitude + phase) images from k-space (step 3-4) and finally converting the phase images to velocity fields (step 4-5) \cite{petersson2010_PCMRIsim,rich2016_PCMRIBayesian,rich2019_4DflowBayesian,zhang2022_4DWSS}. The aim here is to start with a map that models just 4-5, then build 3-5, and progressively attempt to model the whole sequence.



Color Doppler Echocardiography works on the principle of Doppler effect wherein the frequency shift in the transmitted and received echo signals is mapped to the velocity of the moving structure, along the direction of the probe. Unlike MRI, several different approaches to model Echo image generation have been extensively studied in the literature \cite{verweij2014_EchoSim}. Several modeling assumptions which include transducer characteristics \cite{jensen1992_echosyn}, imaging media properties (attenuation, absorption dispersion)\cite{wells1975_Echoabsoprtion,treeby2011_Echoattenuation}, near and far field behaviour \cite{verweij2014_EchoSim} have led to numerous ultrasound image reconstruction techniques. Hence, careful assessment of the considerations proposed by previous studies is crucial in formulating the Echo measurement operator.

\subsubsection*{Task 1.3: Develop a unifying approach to solve the co-registration problem.}

One of the significant challenges in medical imaging is the co-registration problem, i.e., the process of aligning and overlaying the image coordinates of different single-patient modalities. 
Patient motion can be one of the causes for the misalignment; the primary reason is that two acquisitions may have been taken at different times or from different modalities imaging at different planes and angles.
We hypothesize that the co-registration problem can be posed in a unifying manner as part of our overall inverse problem formulation.

Typically, all medical imaging protocols incorporate EKG/ECG gating.
So, we assume that the times are appropriately matched to simplify the problem description.
Let $\boldsymbol{\phi}_i = (g_i,\mathbf{v}_i,\dots)$ be the collection of hemodynamic fields pertaining to modality $i$.
We will follow a hierarchical Bayesian approach in modeling these fields (see Fig.~\ref{fig:coregistration}).
Specifically, assume that these fields follow a global pattern $\boldsymbol{\phi}_0$ with small variations $\delta\boldsymbol{\phi}_i$ for each modality, e.g.,
$
\mathbf{v}_i = \mathbf{v}_0 + \delta \mathbf{v}_i
$
and so on.

\begin{wrapfigure}{r}{.7\textwidth}
    % \vspace*{\fill}
    \centering
    \begin{subfigure}[b]{0.4\textwidth}
    \includegraphics[width=\textwidth]{figs/coregistration-dag.pdf}
    %\caption{Causal graph of the data Hamiltonian.}
    \end{subfigure}
    % \label{fig:suba}} %\par\vfill
    \begin{subfigure}[b]{0.29\textwidth}
    \includegraphics[width=\textwidth,trim={0 0 0 0},clip]{figs/field-transformation.pdf}
    %\caption{Patient movement induces a transformation on the underlying field.}
    \end{subfigure}
    % \label{fig:subb}}
    \caption{Co-registration as a Bayesian hierarchical inference problem.}
    \label{fig:coregistration}
\end{wrapfigure}
We capture the patient's position relative to the imaging plane by a translation, $\boldsymbol{\tau}_i(t)$ in $\mathbb{R}^3$, and a rotation, $\mathbf{U}_i(t)$ in $\text{SO}(3)$.
Let $T_{\boldsymbol{\tau}_i,\mathbf{U}_i}$ denote an operator that acts on the fields on $\boldsymbol{\phi}_i$ and applies the translation $\boldsymbol{\tau}_i$ and the rotation $\mathbf{U}_i$ at each time.
The order does not matter because they commute.
The action of this operator depends on the field. For example,
$$
\small
\left(T_{\boldsymbol{\tau}_i,\mathbf{U}_i}g_i\right)(x,t) = g_i(\mathbf{U}_i^T(t)x -\boldsymbol{\tau}_i(t),t),\;\text{and}\;
\left(T_{\boldsymbol{\tau}_i,\mathbf{U}_i}\mathbf{v}_i\right)(x,t) = \mathbf{U}_i^T(t)\mathbf{v}_i(\mathbf{U}_i^T(t)x - \boldsymbol{\tau}_i(t),t),
$$
are the transformations for geometry (scalar) and velocity (vector), respectively (see Fig.~\ref{fig:coregistration}).
We denote the action of $T_{\boldsymbol{\tau}_i,\mathbf{U}_i}$ on each component of $\boldsymbol{\phi}_i$ by $T_{\boldsymbol{\tau}_i,\mathbf{U}_i}\boldsymbol{\phi}_i$.

We can now define the data Hamiltonian.
Suppose we have measurements at $n_i$ different timesteps $t_{ij}$.
Let $d_{ij}$ be the measurement of modality $i$ at time $t_j$.
Assuming that the measurements are independent conditional on the fields, we have the following additive decomposition:
$$
\small
H(d_i|R_i\boldsymbol{\phi}_i) = H\left(d_{i1}|R_{i,t_{i1}}^{\text{ref}}T_{\boldsymbol{\tau}_i(t_{i1}),\mathbf{U}_i(t_{i1})}\boldsymbol{\phi}_i\right)
+\dots+
H\left(d_{in}|R_{i,t_{in}}^{\text{ref}}T_{\boldsymbol{\tau}_i(t_{in}),\mathbf{U}_i(t_{in})}\boldsymbol{\phi}_i\right)
.
$$
The term inside the sum depends on the measurement noise model.
The simplest choice would be a zero mean multivariate Gaussian noise with covariance matrix $N_i$:
$$
\small
H\left(d_{ij}|R_{i,t_{ij}}^{\text{ref}}T_{\boldsymbol{\tau}_i(t_{ij}),\mathbf{U}_i(t_{ij})}\boldsymbol{\phi}_i\right)
= \frac{1}{2}\left(R_{i,t_{ij}}^{\text{ref}}T_{\boldsymbol{\tau}_i(t_{ij}),\mathbf{U}_i(t_{ij})}\boldsymbol{\phi}_i -d_{ij}\right)^TN_i^{-1}\left(R_{i,t_{ij}}^{\text{ref}}T_{\boldsymbol{\tau}_i(t_{ij}),\mathbf{U}_i(t_{ij})}\boldsymbol{\phi}_i - d_{ij}\right).
$$
The all-modality data Hamiltonian is the sum of these Hamiltonians over $i$.

The co-registration problem requires us to find the global field pattern $\boldsymbol{\phi}_0$, the modality variations, $\delta \boldsymbol{\phi}_i$, and the time-dependent transformations $\boldsymbol{\tau}_i(t)$ and $\mathbf{U}_i(t)$.
For the hemodynamic fields, we will use the physics-informed Hamiltonians of $\emph{Task 1.1}$ on each modality.
We will define a suitable information Hamiltonian for $\boldsymbol{\tau}_i$ and $\mathbf{U}_i$ as part of this task.

\subsection*{Specific Aim 2: Parameterization of time-evolving hemodynamic fields.}

\paragraph{Introduction.}
The objective is to study the properties of different parameterizations of spatiotemporal hemodynamic fields. Typically, in IFT, one parameterizes fields by expanding them in a basis, most commonly piecewise linear functions over voxels or Fourier. One then characterizes the posterior over the parameters instead of the infinite-dimensional measure over the function space; see previous discussion on renormalization groups, Sec.~\ref{sec:IFTreview}. 
Furthermore, in IFT, it is standard to treat time as another spatial variable.
This treatment ignores the Markovian nature of time evolution.
We develop an approach that explicitly accounts for time evolution.

\subsubsection*{Task 2.1: Parameterization of spatial fields.}
We will start by studying the advantages and disadvantages of different parameterizations.
In particular, we are interested in the following: 
1) Assessing in what sense the induced finite dimensional probability measure over the field parameters converges to the IFT function space measure;
2) The smallest possible scale that a given parameterization can represent;
3) Whether there are parameterizations that can be locally adapted to represent sharp features;
4) Convergence properties of the parameterization to fields.

Investigating convergence in probability to the IFT measure may require showing that the measure is well-defined.
In the Gaussian case, this is trivial.
The non-Gaussian case will require more work.
We will explore this by defining the path integral~\cite{ivanov2021notes}, or by constructing
\begin{wrapfigure}{r}{0.45\textwidth}
  \begin{center}
    \includegraphics[width=0.44\textwidth,trim={0 0 0 0},clip]{figs/parameterized-field.pdf}
  \end{center}
  \caption{Field parameterization.}
  \label{fig:parameterized-field}
\end{wrapfigure}
 generalized random fields by L\'evy's continuity theorem~\cite{bierme2017generalized}.
We will also investigate how the parameterization converges to fields in the limit of infinite parameters.
Under the Paley-Zygmund theorem, there are certain conditions for which a random Fourier series almost surely does not converge to any function in $L^p$~\cite{hill2012convergence}.
We must determine if the convergence conditions are satisfied for a given field parameterization.
We will investigate theoretically and numerically the behavior of different field parameterizations, including piecewise linear, Fourier, wavelets, radial basis functions, and DNNs.

We will also study how to enforce boundary conditions parametrically.
For example, we can enforce the no-slip boundary conditions by:
$
\mathbf{v} = g\mathbf{u},
$
where $\mathbf{u}$ is an arbitrarily parametrized vector field.
Notice that $g=0$ at the vessel boundary; thus, $\mathbf{v}$ is also zero.
The pressure is arbitrary up to an additive constant.
The arbitrariness can be removed by parameterizing as $\tau(x,t) = \phi(x,t) - \phi(x',t)$ where $x'$ is an arbitrary spatial point inside the vessel and $\phi$ is an arbitrarily parameterized field.
For initial conditions, see \emph{Task 2.2}.

\subsubsection*{Task 2.2: Parameterization of time evolution.}
We will derive a field parameterization of the fields that enforces initial conditions and the Markovian property (see Fig. \ref{fig:parameterized-field}).
Recall that $\boldsymbol{\phi}$ denotes the tuple of all fields.
Suppose that $\hat{\boldsymbol{\phi}}(x;\boldsymbol{\theta})$ is a specific spatial parameterization of these fields using parameters $\boldsymbol{\theta}$.
The dynamics of the field induce dynamics on the space of parameters $\boldsymbol{\theta}$:
$$
\small
\dot{\boldsymbol{\theta}} = \mathbf{f}(\boldsymbol{\theta},t),
$$
with initial conditions $\theta(0) = \theta_0$.
Here, we have to make $\mathbf{f}$ depend on time because the system we are studying is not closed and follows the cardiac cycle.
Our idea is to parameterize $\mathbf{f}$:
$$
\mathbf{f}(\boldsymbol{\theta},t) \approx \hat{\mathbf{f}}(\boldsymbol{\theta},t;\boldsymbol{\psi}),
$$
with parameters $\boldsymbol{\psi}$, an idea similar to Neural ODEs~\cite{chen2018neural}.
Expanding our work on dynamical systems~\cite{hao2023information}, we can define a probability measure over $\boldsymbol{\theta}$ conditional on the initial conditions.
Specifically, if $\boldsymbol{\theta}_{-0}$ is the function $\boldsymbol{\theta}$ for $t>0$, we define $p(\boldsymbol{\theta}_{-0}|\boldsymbol{\theta}_0)$ through the information Hamiltonian:
$$
\small
H(\boldsymbol{\theta}_{-0}) = -\beta H(\hat{\boldsymbol{\phi}}\left(\cdot;\boldsymbol{\theta})\right) -\beta_{\boldsymbol{\theta}}\int_0^T dt |\dot{\boldsymbol{\theta}}-\hat{\mathbf{f}}(\boldsymbol{\theta},t;\boldsymbol{\psi})|^2,
$$
where $T$ is the terminal time, and $H(\hat{\boldsymbol{\phi}}\left(\cdot;\boldsymbol{\theta})\right)$ is the information Hamiltonian of $\boldsymbol{\phi}$ evaluated at the parameterization.
The problem now becomes to infer $\boldsymbol{\theta}$ and the dynamic parameters $\psi$.
We can parameterize $\boldsymbol{\theta}$ so that it satisfies the initial conditions, e.g., if $\mathbf{h}(t)$ is an arbitrary parameterization:
$$
\small
\boldsymbol{\theta}(t;\boldsymbol{\theta}_0) = \boldsymbol{\theta}_0 + t \mathbf{h}(t).
$$

\subsection*{Specific Aim 3: Numerical algorithms for characterizing joint posterior.}
\paragraph{Introduction.}
The objective is to develop scalable numerical schemes to characterize the joint posterior over hemodynamic fields and physical parameters. There are two sources of computational complexity: the data dimension and the evaluation of Hamiltonians. We seek algorithms that can process data and physics-enforcing equations in batches and parallel with minimal communication. We will focus on sampling approaches that are variants of nested, non-convergent, persistent SGLD, and variational inference. 

\subsubsection*{Task 3.1: Advanced stochastic gradient Langevin dynamics}
We will derive scalable sampling algorithms to characterize the joint posterior of fields and parameters.
Sampling approaches are necessary because they do not introduce biases.
We will exploit the data parallelism framework to scale to large data sets and spatiotemporal domains.
Let $\lambda$ denote all parameters on which the partition functions of the various Hamiltonians depend.
The SGLD sampling algorithm we developed in~\cite{alberts2023physics} requires a noisy estimate of the gradient of:
$$
\small
H(\lambda|d) = -\log p(\lambda | d) = -\log \int \mathcal{D}\boldsymbol{\phi} \: p(\boldsymbol{\phi},\lambda |d).
$$
While this \emph{parameter information Hamiltonian} is intractable due to the path integral, its gradient $\nabla_{\lambda}H(\lambda |d)$ can be expressed through field expectations.
Hence, one can construct a noisy estimate that avoids evaluating the path integral.
This estimate is sufficient to perform SGLD.

The main drawback of SGLD is that the posterior samples mix very slowly.
One attempt to improve SGLD is using \emph{preconditioned} SGLD~\cite{li2016preconditioned}, which uses a preconditioning matrix in the update for better convergence behavior.
The choice of the preconditioning matrix is up to the user, although matrices that use curvature information perform better.
%(see Fig. \ref{fig:preconditioned}).

% \begin{wrapfigure}{r}{0.45\textwidth}
%   \begin{center}
%     \includegraphics[width=0.44\textwidth,trim={0 0.1cm 0 0.4cm},clip]{figs/preconditioned-sgld.pdf}
%   \end{center}
%   \caption{Sampling the IFT posterior can be made more efficient by using a preconditioning matrix.}
%   \label{fig:preconditioned}
% \end{wrapfigure}

We can also express the Hessian of the parameter information Hamiltonian, $\nabla_{\lambda}^2H(\lambda| d)$, using expectations.
Doing so will provide an ideal preconditioning matrix since the Hessian contains information about the curvature of the landscape.
Some approaches continuously use Hessian information; see~\cite{simsekli2016stochastic}.
The reference shows that one needs a correction requiring third-order derivatives to have a convergent scheme, see~\cite{girolami2011riemann}.
We will investigate the construction of the Hessian and third-order correction for use in improving SGLD and compare it to other common preconditioning choices such as RMSprop~\cite{hinton2012neural}, the expected Fisher information matrix~\cite{patterson2013stochastic} or batch normalization~\cite{lange2023batch}.


\subsubsection*{Task 3.2: Amortized, variational inference}
\begin{wrapfigure}{r}{0.43\textwidth}
  \begin{center}
    \includegraphics[width=0.42\textwidth,trim={0 0 0 0},clip]{figs/amortization.pdf}
  \end{center}
  \caption{Regular vs. amortized variational inference.}
  \label{fig:amortization}
\end{wrapfigure}
Sampling approaches take several iterations to converge.
Variational inference (VI) poses an optimization problem that fits the joint posterior to a parameterized distribution family.
Traditional VI cannot be used in IFT when the partition function depends on parameters.
In~\cite{ha2018world}, we developed a nested version of VI for IFT and observed orders of magnitude faster performance than sampling.
Still, VI for the hemodynamic field reconstruction will not be computationally cheap.
To accelerate inference, we propose to exploit the idea of amortization, see our relevant work in~\cite{KARUMURI2023111851}.
In amortized VI, one poses a problem that learns the map from data to the posterior (see Fig.~\ref{fig:amortization}).
When new data become available, this map provides a quick estimate of the posterior, which could be taken as is or used as a starting point for sampling or VI.
As in the previous task, our proposal is scalable and can use the data parallelism paradigm.

We need two guide distributions, $q_{\zeta}\left(\boldsymbol{\theta}\right)$ and $q_{\eta}\left(\lambda\right)$.
To learn the guide parameters $\zeta$ and $\eta$, we maximize the evidence lower bound (ELBO)~\cite{2013arXiv1312.6114K},
\begin{align*}
\small
    \operatorname{ELBO}
    \left(
    \zeta, \eta \middle\vert d 
    \right)
    =
    \mathbb{E}_{q_{\zeta}(\boldsymbol{\theta}), q_{\eta}(\lambda)}
    \left[
    -H(d|R\hat{\boldsymbol{\phi}}\left(\cdot;\boldsymbol{\theta}\right))
    - 
    H(\hat{\boldsymbol{\phi}}\left(\cdot;\boldsymbol{\theta}\right), \lambda)
    - \log Z(\lambda)
    \right]
    +
    \mathbb{H}\left[
    q_{\zeta}(\boldsymbol{\theta})
    \right]
    +
    \mathbb{H}\left[
    q_{\eta}(\lambda)
    \right],
\end{align*}
where $\mathbb{H}$ is the entropy.
The optimal guides depend on the measurement data $d$.
To achieve real-time inference, we will explore the idea of amortization.

In amortization, we learn data-parameterized guides $q_{\zeta(d)}(\boldsymbol{\theta})$ and $q_{\eta(d)}(\lambda)$ by maximizing the expectation of the ELBO running over all possible measurement data:
\begin{align*}
\small
    \operatorname{AELBO}
    \left[
    \zeta, \eta\middle\vert p(d)
    \right]
    =
    \mathbb{E}_{p(d)}
    \left[
     \operatorname{ELBO}
    \left(
    \zeta(d), 
    \eta(d)
    \middle\vert d 
    \right)
    \right].
\end{align*}
The amortized guides are $q_{\operatorname{NN}(d;\psi_{\zeta})}(\boldsymbol{\theta})$ and
$q_{\operatorname{NN}(d;\psi_{\eta})}(\lambda)$.
The expectation is over all data compatible with the models, but in practice, we can use a set of synthetic datasets.

\subsection*{Specific Aim 4: Verification, validation, and demonstration.}
\paragraph{Introduction.}
The objective is to verify, validate, and demonstrate the proposed methodology. 
We will use 4D Flow MRI and Echo data generated through simulation for verification. 
Validation will involve in vitro experiments across three modalities, benchmarked against PIV/PTV experiments conducted on identical phantoms. 
\begin{comment} CFD simulation results from previous publications/publically available datasets would also be used. (NOTE: Uncomment this line only after finding some relevant literature/datasets) 
\end{comment}
For demonstration, we will apply the methodology to clinical datasets covering a range of heart conditions. 
We will compare our results to those of alternative methods in each phase. 
Further information can be found in Table \ref{table:datasets}.

\begin{table*}[!h]
\scriptsize
    \centering
    \caption{Outline of verification, validation, and demonstration process.}
    \label{table:datasets}
    \begin{tabularx}{\textwidth}{c | X | X | X}
        \toprule
        \textbf{Step} & \textbf{Datasets} & \textbf{Metrics} &  \textbf{Methods Compared} \\
        \midrule
        \multirow{8}{*}[2ex]{\centering Verification} & 
        \textbf{Data}:
        \begin{itemize}[left=0em]
            \item Synthetic 4D Flow
            \item Echo
        \end{itemize}
        \textbf{Groundtruth}:
        \begin{itemize}[left=0em]
            \item Analytical flow solution
        \end{itemize} & 
        \begin{itemize}[left=0em]
            \item RMS error distribution
            \item Bland–Altman analysis
        \end{itemize} & 
        \multirow{32}{*}[16ex]{%
            \begin{minipage}{.25\textwidth}
                \begin{itemize}[left=0em]
                    \item PC-MRI
                        \begin{itemize}
                            \item Loecher et al. \cite{loecher2018_PCMRI}
                        \end{itemize}
                    \item 4D Flow MRI
                        \begin{itemize}
                            \item Zhang et al. \cite{zhang2022_4Dsparserep, zhang2019_4Dpressure, zhang2022_4DWSS}
                            \item Ferdian et al. \cite{Ferdien2020_4DFlowNN}
                        \end{itemize}
                    \item Echo
                        \begin{itemize}
                            \item DoVeR \cite{brett2020_echo}
                            \item iVFM \cite{assi2017_EchoIVFM}
                        \end{itemize}
                    \item Boundary segmentation
                        \begin{itemize}
                            \item Manual segmentation
                            \item Time-of-Flight MRA (semi-automated using ITK-SNAP \cite{yushkevich2006_ITKSNAP})
                            \item SDM velocity \cite{rothenberger2023_4DflowSDM}
                        \end{itemize}
                    \item PINNs
                        \begin{itemize}
                            \item PINNs~\cite{cai2021physics}
                            \item Bayesian PINNs~\cite{yang2021b}
                        \end{itemize}
                \end{itemize}
            \end{minipage}
        } \\
        \cmidrule{1-3}
        \multirow{8}{*}[2ex]{\centering Validation} & 
        \textbf{Data}:
        \begin{itemize}[left=0em]
            \item In Vitro 4D Flow: Poiseuille flow
            \item Patient-specific aneurysms
        \end{itemize}
        \textbf{Groundtruth}:
        \begin{itemize}[left=0em]
            \item 3D-PIV/PTV measurements
        \end{itemize} & 
        \begin{itemize}[left=0em]
            \item 2D contours of the field
            \item 3D streamlines, pathlines
            \item Bland–Altman analysis
            \item Pearson's correlation across modalities
        \end{itemize} & \\
        \cmidrule{1-3}
        \multirow{8}{*}[2ex]{\centering Demonstration} & 
        \textbf{Data (clinical)}:
        \begin{itemize}[left=0em]
            \item Mouse Left ventricle (LV)
            \item Human Right Ventricle (RV)
            \item Aortic flow
        \end{itemize}
        \textbf{Groundtruth}: NA & 
        \begin{itemize}[left=0em]
            \item q-q plot
            \item AIC, BIC
            \item 3D streamlines, pathlines
            \item 2D field contours
            \item Pearson's correlation across modalities
        \end{itemize} & \\
        \bottomrule
    \end{tabularx}
\end{table*}



\subsubsection*{Task 4.1: Verify methodology using synthetic experimental data.}
We will generate synthetic medical images based on known analytical flow fields over simple geometries for each modality. 
To accomplish this, we will use previously developed in-house codes for MRI \cite{zhang2021_4Dphaseunwrap,zhang2022_4Dsparserep} and publicly available resources for Echo \cite{jensen1992_echosyn, jensen1997_echosyn, sun2022_echodata}.

After generating the synthetic images, we will systematically study the reconstruction performance of our method for diverse input conditions. 
We will examine how well the methodology performs under various flow patterns, noise levels, and other image artifacts often encountered in real-world applications.
In addition, we will also conduct a convergence analysis by studying the effects of initial conditions and parameter initialization on the quality of the reconstructed fields. 
Through this analysis, we aim to understand how sensitive the method is to these factors and whether it reliably converges to an accurate solution.

For the quantitative assessment, we will use metrics like RMS error distribution to measure the accuracy of the reconstructed flow fields. 
We will use Bland–Altman analysis to compare the method's output with the known ground truth, to assess its limits and potential systematic biases.

\subsubsection*{Task 4.2: Validate methodology using in vitro experimental data.}

We will create a workflow for in vitro experiments focusing on patient-specific aneurysm geometries. 
MRI experiments will take place at the on-campus MRI facility, while Echo scans will be collected using Dr. Vlachos' lab equipment. 
Our team will apply its PIV/PTV expertise to conduct volumetric flow experiments on identical flow phantoms~\cite{Brindise2019Multi-modality}.
The PIV/PTV measurements will serve as the ground truth for comparison against results obtained through our method.

Given the constraints each imaging modality imposes on the setup's materials choice, we will construct a compatible flow loop across all modalities. This ensures the setup remains constant while switching between MRI, Echo, and PIV/PTV experiments. 
Finally, we will pit the results from our proposed IFT reconstruction method against those obtained from PIV/PTV. 

For qualitative and quantitative evaluation, we will use a range of methods. 
These will include 2D field contours and 3D visualizations like streamlines and pathlines. 
Bland–Altman analysis will help assess the limits of agreement between the methods, and Pearson's correlation will provide a measure of consistency across different imaging modalities.
This comparison will serve as the basis for validating the accuracy and reliability of our method in a controlled experimental setting.

\begin{figure}%[htbp]
    \centering
    \includegraphics[width=\textwidth]{figs/MRI_invitro.pdf}
    \caption{ In vitro experiments on patient-specific aneurysm geometry; (i) a pdms model created for PTV experiments; (ii) Comparison of reconstructions from in vivo 4D Flow MRI and in vitro PTV~\cite{brindise_multi-modality_2019,zhang2019_4Dpressure}.}
    \label{fig:invitro_MRI}
\end{figure}

\subsubsection*{Task 4.3: Demonstrate methodology on clinical data}
We will apply the methodology to clinical data to reconstruct hemodynamic fields. 
The clinical datasets include small animal models like a healthy mouse left ventricle (50 frames per cycle) ~\cite{brett2020_echo}, and human subjects, covering areas such as the right ventricle and aortic (30 patients total - 12 bicuspid, 12 tricuspid+aneurysm, 6 healthy) flow~\cite{zhang2021_4Dphaseunwrap,zhang2022_4DWSS}. 
Through existing collaborations with Indiana University, Children's National Hospital, and North Western University, we will gain access to retrospective data on specific cardiac diseases, including Single Ventricle defects and Tetralogy of Fallot (20 patients) \cite{brett2023_EchorTOF, brett2022_Echo4D}. Access to this diverse set of clinical data will allow us to assess the predictive and risk assessment capabilities of our methodology.

For evaluation, the focus will be on multiple metrics. 
Q-q plots will serve to analyze how closely the data follows a particular distribution, thus validating underlying statistical assumptions.
Akaike information criterion (AIC) \cite{sakamoto1986akaike} and Bayesian information criterion (BIC) \cite{vrieze2012model} will be used to compare different models based on their goodness-of-fit and complexity, offering insights into the most suitable models for the clinical data at hand. 
These metrics will help us understand the methodology's performance in real-world clinical scenarios.

\begin{figure}%[htbp]
    \centering
    \includegraphics[width=\textwidth]{figs/dover_invivo_result.pdf}
    \caption{Vorticity field reconstructions for a mouse LV Echo data: Method comparison~\cite{brett2020_echo}.}
    \label{fig:invivo_echo_result}
\end{figure}

\section{Broader Impacts\label{sec:BI}}

% The Project Description also must contain, as a separate section within the narrative, a section labeled “Broader Impacts”. This section should provide a discussion of the broader impacts of the proposed activities. Broader impacts may be accomplished through the research itself, through the activities that are directly related to specific research projects, or through activities that are supported by, but are complementary to the project. NSF values the advancement of scientific knowledge and activities that contribute to the achievement of societally relevant outcomes. Such outcomes include, but are not limited to: full participation of women, persons with disabilities, and underrepresented minorities in science, technology, engineering, and mathematics (STEM); improved STEM education and educator development at any level; increased public scientific literacy and public engagement with science and technology; improved well-being of individuals in society; development of a diverse, globally competitive STEM workforce; increased partnerships between academia, industry, and others; improved national security; increased economic competitiveness of the U.S.; use of science and technology to inform public policy; and enhanced infrastructure for research and education. These examples of societally relevant outcomes should not be considered either comprehensive or prescriptive. Proposers may include appropriate outcomes not covered by these examples.


We have planned the following specific broader impact tasks:
\begin{itemize}[leftmargin=*]
\item \textbf{Broader Impact Task 1: Publish open source software}
We will develop open-source software aligned with our Bayesian methodologies for cardiovascular hemodynamics. 
This software will be disseminated under the MIT license, ensuring wide accessibility and adaptability. 
Additionally, all code related to our research publications will be made freely available to the public. 
This initiative aims to empower other researchers and domain experts by providing them with the computational tools to implement and explore IFT and other statistical methodologies that we develop in medical imaging.

\item \textbf{Broader Impact Task 2: Publish verification and validation benchmark datasets}
We will release cardiovascular hemodynamics verification and validation benchmark datasets along with data generation codes. 
These datasets will aim to test and validate Bayesian methodologies and other statistical approaches. 
The datasets will be hosted on the Purdue University Research Repository (PURR) with DOIs for citation. 
They will be disseminated under a permissive license to ensure accessibility. Documentation will accompany the datasets to outline the experimental setup, data collection methods, and pre-processing steps. 
This ensures reproducibility and facilitates use. 
The goal is to provide a framework for researchers to compare and validate methodologies, accelerating progress in the field and fostering collaboration for future studies.

\item \textbf{Broader Impact Task 3: Engage undergraduate researchers}
We will involve one undergraduate student per year in our project through Purdue University's Summer Undergraduate Research Fellowship (SURF) program, a 10-week paid research experience.
Students will work alongside faculty and graduate mentors to develop interactive GUI tools that demonstrate our Bayesian methodologies for cardiovascular hemodynamics.
This initiative builds on our past success with the SURF program, where we created R2Roptimizer, a tool for optimizing graphene growth processes~\cite{26898}. 
By engaging undergraduates in this manner, we aim to enrich their educational experience while broadening the reach of our research.

\item \textbf{Broader Impact Task 4: Disseminate research material through undergraduate and graduate courses}
We will develop Jupyter notebook activities which will be disseminated in two undergraduate classes (ME 297 Introduction to Data Science for Mechanical Engineers, and ME 309 Introduction to Fluid Mechanics) and one graduate course ME 539 Introduction to Scientific Computing.
Bilionis has already created dozens of such activities which he has compiled in the form of two online textbooks~\cite{ME297-textbook,ME539-textbook}.

\item \textbf{Broader Impact Task 5: Organize conference mini-symposia}
We have already planned to organize a minisymposium at the Society for Industrial and Applied Mathematics (SIAM) Conference on Uncertainty Quantification (UQ) between February 27 -- March 1, 2024.
The goal of the minisymposium is to provide a platform for researchers working with IFT to share their advancements and engage in meaningful dialogues.
We aim to promote the significance of this emerging topic to the wider community.
We will organize similar minisymposia in SIAM Computational Science and Engineering (CSE) \& UQ conferences in subsequent years, and American Physics Society (APS) Division of Fluid Dynamics (DFD) and in years 2 and 3 on one of the Society for Cardiovascular Magnetic Resonance (CMR), International Society for Magnetic Resonance in Medicine (ISMRM), Society for Magnetic Resonance Angiography (SMRA), or American Society of Echocardiography (ASE).

\item \textbf{Broader Impact Task 6: Present research results to clinical practitioners}
All publishable results and findings will be presented to the clinical audience through our participation in the American Heart Association (AHA), International Symposium for Magnetic Resonance Imaging (ISMRM) annual conferences, or American Society of Echocardiography (ASE).
\end{itemize}

\section{Results of Prior NSF Support}
% ILIAS MUST COLLECT INFO FROM PAVLO

% The purpose of this section is to assist reviewers in assessing the quality of prior work conducted with prior or current NSF funding. If any PI or co-PI identified on the proposal has received prior NSF support including:

% an award with an end date in the past five years; or

% any current funding, including any no cost extensions

% Information on the award is required for each PI and co-PI, regardless of whether the support was directly related to the proposal or not. In cases where the PI or any co-PI has received more than one award (excluding amendments to existing awards), they need only report on the one award that is most closely related to the proposal. Support means salary support, as well as any other funding awarded by NSF, including research, Graduate Research Fellowship, Major Research Instrumentation, conference, equipment, travel, and center awards, etc.

% The following information must be provided:

% (a) the NSF award number, amount and period of support;

% (b) the title of the project;

% (c) a summary of the results of the completed work, including accomplishments, supported by the award. The results must be separately described under two distinct headings: Intellectual Merit and Broader Impacts;

% (d) a listing of the publications resulting from the NSF award (a complete bibliographic citation for each publication must be provided either in this section or in the References Cited section of the proposal); if none, state “No publications were produced under this award.”

% (e) evidence of research products and their availability, including, but not limited to: data, publications, samples, physical collections, software, and models, as described in any Data Management Plan; and

% (f) if the proposal is for renewed support, a description of the relation of the completed work to the proposed work.

% If the project was recently awarded and therefore no new results exist, describe the major goals and broader impacts of the project. Note that the proposal may contain up to five pages to describe the results. Results may be summarized in fewer than five pages, which would give the balance of the 15 pages for the Project Description.


\textbf{I. Bilionis} most relevant prior NSF support. \textbf{Award number:} 1922316-DMREF. \textbf{Amount: } \$1,738,752. \textbf{Period of support:} 10/1/2019 -- present. \textbf{Title:} Discovery of high-temperature, oxidation-resistant, complex, concentrated alloys via data science-driven multi-resolution experiments and simulations. \textbf{Summary:} The project seeks to use a combination of machine learning, physics-based modeling, and experiments to discover Al-containing complex concentrated alloys with high hardness. In addition, the team seeks to advance the physics-based modeling of RCCAs, machine-learning modeling in the context of materials science, and synthesis of RCCAs. \textbf{Intellectual merit:} Machine learning models that predict strength and corrosion properties of RCCAs. Algorithms for the sequential design of experiments of such alloys. Experimental properties of tens of RCCAs.
\textbf{Broader Impacts:} The team has created eight Nano-HUB simulation tools that enable the exploration of RCCAs properties~\cite{35885,31736,33036,33792,36022,36148,35881}. Four graduate students, one undergraduate student, and two postdocs have been trained through this grant. The results have been disseminated in numerous conference presentations and invited talks. \textbf{Publications:} \cite{KARUMURI2023111851, FARACHE2022111386,10.1063/5.0076584,strachan2021,McClure2021} 


\textbf{P. Vlachos} most relevant prior NSF support. \textbf{Award number:} 1805817-CBET. \textbf{Amount: }\$359,655. \textbf{Period of support:} 9/1/2021 -- present. \textbf{Title:} Miscibility-immiscibility conundrum in air-liquid-vapor flow modeling: Bridging the gap by using the phase-field method. \textbf{Summary:} This project developed advanced multi-physics and multi-scale models for the inception and development of cavitation in non-equilibrium systems. \textbf{Intellectual merit:} A new theoretical model of cavitation in porous viscoelastic media was developed by modifying the classical Raleigh-Plesset equation. The model results in a single simplified ODE that predicts the bubble dynamics subject to the heterogeneity of the porous domain. \textbf{Broader Impacts:} This model enables novel investigation of cavitation in biological systems and especially the effect of cavitation on brain tissue during traumatic brain injury or cavitation-mediated drug delivery. \textbf{Publications:} \cite{10.1063/1.5109889, HU2022114526, 10.1063/5.0021697, 10.1093/pnasnexus/pgac150}.

%\section{Management Plan and Timeline\label{sec:personnel}}

%Bilionis is an expert on uncertainty quantification ...
%Vlachos is an expert on experimental flow measurements ...
%Bilionis and Vlachos have been collaborating since 2019 on several projects funded by Eli Lilly.
%Their common work has focused on the reconstruction of flow fields using particle image velocimetry (PIV).
%They have published together [x].
%They have invented a PIV technology that has been disclosed to Eli Lilly but is pending clearance for publication.
%As part of their common Eli Lilly-funded project, they have been meeting at least twice a week since 2019.
%They have formed a joined subgroup of six graduate students and postdocs who meet on a weekly basis to discuss flow field reconstruction in biomedical problems.
%They are currently co-advising two graduate students.

%In this proposed project, Bilionis and Vlachos will co-advise a postdoctoral researcher, see attached \textit{Mentoring Plan}.
%Bilionis will be responsible for the management of the project, and reporting activities.
%On the technical side, Bilionis will lead Task 1.1, and specific aims 2 \& 3.
%Vlachos will lead Tasks 1.2-1.4 and specific aim 4.